% GENERATED FILE. Edit canonical lessons, then run npm run generate:books.
% canonical-source-hash: fnv1a64:210a97e034ab60a2
% canonical-lessons: UR-C100-chauraha, UR-C100-aspas, UR-C100-sath, UR-C100-bazar, UR-C100-thana

\chapter{Crossroads, Around, With, Market, Police station}
\label{ch:ur-crossroads-around-with-market-police-station}
\hlchaptermodality{100}

\begin{chapteropening}
\textbf{By the end of this chapter:} \emph{I can say a crossroads, around, with, a market and a police station.}

Complete the last lesson of chapter 100: I can say a crossroads, around, with, a market and a police station.
\end{chapteropening}

\section[\texorpdfstring{\ur{چوراہا} (chaurāhā)}{چوراہا (chaurāhā)}]{\ur{چوراہا} (chaurāhā) — a crossroads}

\label{lesson:UR-C100-chauraha}



\begin{quote}
Before the new one: say the Urdu for a lane, a narrow street, then the Urdu for a turn, a bend.
\end{quote}

\subsection*{You'll want to know: \ur{چوراہا}}
\textbf{\ur{چوراہا}} — \emph{chaurāhā} — \textquotedblleft{}a crossroads\textquotedblright{}.

\begin{cousinweb}
\textbf{\ur{چوراہا}} (\emph{chaurāhā}) is usually explained as \emph{chau-}, \textquotedblleft{}four\textquotedblright{}, with Persian \emph{rāh}, \textquotedblleft{}road\textquotedblright{}. CDIAL does not list it; it lists \emph{cauhaṭṭā}, \textquotedblleft{}a crossing of market streets\textquotedblright{}, from \emph{catur-}, \textquotedblleft{}four\textquotedblright{} (CDIAL 4626).
\end{cousinweb}

\begin{grammarlens}[title={What you've built}]
One more word for this chapter.
\end{grammarlens}

\subsection*{Your turn}
\begin{itemize}
\raggedright
  \item \emph{Say it:} \emph{chaurāhā}
  \item \emph{Say it:} \emph{chaurāhā}, once more
  \item \emph{Say it:} say \emph{chaurāhā}
  \item \emph{Recall:} say the Urdu for a lane, a narrow street, then the Urdu for a turn, a bend, then say \emph{chaurāhā} again
\end{itemize}

\begin{tcolorbox}[breakable,colback=teal!4,colframe=teal!35!black,title={Before you move on}]
What does \ur{چوراہا} mean? (\textquotedblleft{}A crossroads\textquotedblright{}.) Say it once more.
\end{tcolorbox}
\section[\texorpdfstring{\ur{آس} \ur{پاس} (ās pās)}{آس پاس (ās pās)}]{\ur{آس} \ur{پاس} (ās pās) — around, nearby}

\label{lesson:UR-C100-aspas}



\begin{quote}
Before the new one: say the Urdu for a turn, a bend, then the Urdu for a crossroads.
\end{quote}

\subsection*{You'll want to know: \ur{آس} \ur{پاس}}
\textbf{\ur{آس} \ur{پاس}} — \emph{ās pās} — \textquotedblleft{}around, nearby\textquotedblright{}.

\begin{cousinweb}
\textbf{\ur{آس} \ur{پاس}} (\emph{ās pās}) is built on \emph{pās}, \textquotedblleft{}near\textquotedblright{}, from Sanskrit \emph{pārśve}, \textquotedblleft{}near\textquotedblright{}, a form of \emph{pārśva}, \textquotedblleft{}side\textquotedblright{} (CDIAL 8118). \emph{ās} is a rhyming partner put in front of it; its own source is not settled.
\end{cousinweb}

\begin{grammarlens}[title={What you've built}]
One more word for this chapter.
\end{grammarlens}

\subsection*{Your turn}
\begin{itemize}
\raggedright
  \item \emph{Say it:} \emph{ās pās}
  \item \emph{Say it:} \emph{ās pās}, once more
  \item \emph{Say it:} say \emph{ās pās}
  \item \emph{Recall:} say the Urdu for a turn, a bend, then the Urdu for a crossroads, then say \emph{ās pās} again
\end{itemize}

\begin{tcolorbox}[breakable,colback=teal!4,colframe=teal!35!black,title={Before you move on}]
What does \ur{آس} \ur{پاس} mean? (\textquotedblleft{}Around, nearby\textquotedblright{}.) Say it once more.
\end{tcolorbox}
\section[\texorpdfstring{\ur{ساتھ} (sāth)}{ساتھ (sāth)}]{\ur{ساتھ} (sāth) — with, together}

\label{lesson:UR-C100-sath}



\begin{quote}
Before the new one: say the Urdu for a crossroads, then the Urdu for around.
\end{quote}

\subsection*{You'll want to know: \ur{ساتھ}}
\textbf{\ur{ساتھ}} — \emph{sāth} — \textquotedblleft{}with, together\textquotedblright{}.

\begin{cousinweb}
\textbf{\ur{ساتھ}} (\emph{sāth}) is inherited. It comes from Sanskrit \emph{sārthena}, \textquotedblleft{}in company with\textquotedblright{}, a form of \emph{sārtha}, \textquotedblleft{}caravan, company\textquotedblright{} (CDIAL 13364).
\end{cousinweb}

\begin{grammarlens}[title={What you've built}]
One more word for this chapter.
\end{grammarlens}

\subsection*{Your turn}
\begin{itemize}
\raggedright
  \item \emph{Say it:} \emph{sāth}
  \item \emph{Say it:} \emph{sāth}, once more
  \item \emph{Say it:} say \emph{sāth}
  \item \emph{Recall:} say the Urdu for a crossroads, then the Urdu for around, then say \emph{sāth} again
\end{itemize}

\begin{tcolorbox}[breakable,colback=teal!4,colframe=teal!35!black,title={Before you move on}]
What does \ur{ساتھ} mean? (\textquotedblleft{}With, together\textquotedblright{}.) Say it once more.
\end{tcolorbox}
\section[\texorpdfstring{\ur{بازار} (bāzār)}{بازار (bāzār)}]{\ur{بازار} (bāzār) — a market}

\label{lesson:UR-C100-bazar}



\begin{quote}
Before the new one: say the Urdu for around, then the Urdu for with.

One ending turns a place into a describing word: say \textbf{\ur{لاہور}}, then \textbf{\ur{لاہوری}}, from Lahore.
\end{quote}

\subsection*{You'll want to know: \ur{بازار}}
\textbf{\ur{بازار}} — \emph{bāzār} — \textquotedblleft{}a market\textquotedblright{}.

\begin{cousinweb}
\textbf{\ur{بازار}} (\emph{bāzār}) is a loan, usually traced to Persian \emph{bāzār}, \textquotedblleft{}market\textquotedblright{}.
\end{cousinweb}

\begin{grammarlens}[title={What you've built}]
One more word for this chapter.
\end{grammarlens}

\subsection*{Your turn}
\begin{itemize}
\raggedright
  \item \emph{Say it:} \emph{bāzār}
  \item \emph{Say it:} \emph{bāzār}, once more
  \item \emph{Say it:} say \emph{bāzār}
  \item \emph{Recall:} say the Urdu for around, then the Urdu for with, then say \emph{bāzār} again
\end{itemize}

\begin{tcolorbox}[breakable,colback=teal!4,colframe=teal!35!black,title={Before you move on}]
What does \ur{بازار} mean? (\textquotedblleft{}A market\textquotedblright{}.) Say it once more.
\end{tcolorbox}
\section[\texorpdfstring{\ur{تھانہ} (thānā)}{تھانہ (thānā)}]{\ur{تھانہ} (thānā) — a police station}

\label{lesson:UR-C100-thana}



\begin{quote}
Before the new one: say the Urdu for with, then the Urdu for a market.
\end{quote}

\subsection*{You'll want to know: \ur{تھانہ}}
\textbf{\ur{تھانہ}} — \emph{thānā} — \textquotedblleft{}a police station\textquotedblright{}.

\begin{cousinweb}
\textbf{\ur{تھانہ}} (\emph{thānā}) is a tadbhava, a word inherited from Sanskrit and reshaped by sound change. It comes from \emph{sthāna}, \textquotedblleft{}firm stance; place of abode\textquotedblright{} (CDIAL 13753). CDIAL glosses the Hindi word \textquotedblleft{}police station\textquotedblright{}.
\end{cousinweb}

\begin{grammarlens}[title={What you've built}]
One more word for this chapter.
\end{grammarlens}

\subsection*{Your turn}
\begin{itemize}
\raggedright
  \item \emph{Say it:} \emph{thānā}
  \item \emph{Say it:} \emph{thānā}, once more
  \item \emph{Say it:} say \emph{thānā}
  \item \emph{Recall:} say the Urdu for with, then the Urdu for a market, then say \emph{thānā} again
\end{itemize}

\begin{tcolorbox}[breakable,colback=teal!4,colframe=teal!35!black,title={Before you move on}]
What does \ur{تھانہ} mean? (\textquotedblleft{}A police station\textquotedblright{}.) Say it once more.
\end{tcolorbox}
